% 本文件由 LaTeX 排版 Agent 根据有效 translation.json 自动生成。
% author=Augustine of Hippo; work_id=1314; title=论隐修士的工作

\chapter{\WorkTitle{论隐修士的工作}}

% structure: p0001 source=h1 target=omitted reason=page-title-already-rendered-by-parent

出自\WorkTitle{订正录}，卷二、第二十一章：撰写\WorkTitle{论隐修士的工作}一书，是出于此迫不得已的原因。当迦太基开始有隐修院时，有些隐修士服从宗徒，亲手劳作维持生活；但另一些却愿靠信友的献仪过活，不做任何可获取或供给生活必需品的工作，反而以为并夸耀自己更实践了福音的诫命，因为主说：\BibleQuote{你们仰望天空的飞鸟和田间的百合花} \BibleRef{玛 6:26}。因此，甚至在目的较低但仍热心虔敬的平信徒中，也开始发生骚乱争执，使教会受到困扰，有人支持这一方，有人支持另一方。此外，有些不工作的人留长发。于是，谴责者与辩护者之间的争论因党派情绪而加剧。为此，同城教会可敬的长老奥勒留主教命我对此事略作著述，我便遵命。本书以\Quote{谨遵你之命，圣洁的奥勒留弟兄}开篇。\par

本书在\WorkTitle{订正录}中置于\WorkTitle{论婚姻之善}之后，后者属于公元401年。\par

1. 你的命令，圣洁的奥勒留弟兄，我应当以更大的热诚来服从，因为我看清了是谁借着你发出这道命令。因为我们的主耶稣基督住在你的内心深处，激发你父性和兄弟般的爱德关怀：我们那些忽视服从有福的宗徒保禄的、被称为隐修士的儿子和弟兄们，是否该允许他们享有那种自由？保禄宗徒说：\BibleQuote{谁若不愿工作，就不应当吃饭}。祂借着你的意志和言语来完成祂的工作，通过你命令我，要我对此事写一些东西给你。因此，愿祂也与我同在，使我如此服从，以至于从祂恩赐的富有成效的劳作之益处中，明白我确实是在服从祂。\par

2. 首先，我们该看看那些主张不工作的修道人说了什么；然后，如果我们发现他们想法不对，该说什么来纠正他们。他们说道：\Quote{宗徒所给的命令，并非由于这种农夫或工匠所从事的体力工作，当他说：\BibleQuote{谁若不愿工作，就不应当吃饭}。因为他不能与福音相悖，主亲自说：\BibleQuote{因此我告诉你们：不要为你们的生命忧虑吃什么，也不要为你们的身体忧虑穿什么。生命不是贵于食物，身体不是贵于衣服吗？你们仰观天空的飞鸟，它们不播种、不收割，也不在仓里积存；你们的天父尚且养活它们；你们不是比它们更贵重吗？你们中谁能用思虑使自己的身量增加一肘呢？关于衣服，你们又为什么忧虑？你们观察田间的百合花怎样生长；它们不劳作，也不纺织；但我告诉你们：连撒罗满在他极荣华时所穿戴的，也不如这花中的一朵。田野的草今天还在，明天就投入炉中，天主尚且这样装饰它，何况你们呢！小信德的人啊！所以你们不要忧虑说：我们吃什么，喝什么，穿什么？因为这都是外邦人所寻求的；你们的天父原知道你们需要这一切。你们先该寻求天主的国和它的义德，这一切自会加给你们。所以你们不要为明天忧虑，因为明天有明天的忧虑：一天的苦足够一天的了。}}。看，他们说，主命令我们不要为饮食衣服忧虑；那么宗徒怎能与主相反，教导我们该如此忧虑吃什么、喝什么、穿什么，甚至用工匠的技艺、操心、劳作来加重我们的负担呢？因此，当他说：\BibleQuote{谁若不愿工作，就不应当吃饭}时，我们必须理解为属灵的工作；关于这工作，他在另一处说：\BibleQuote{按主所赐的，我栽种了，阿颇罗浇灌了，然而使之生长的却是天主}。稍后又说：\BibleQuote{各人要按自己的劳苦领取自己的报酬。我们是天主的同工，你们是天主的田地，天主的建筑物。按照天主所赐我的恩宠，我像一位精明的建筑师，奠立了根基}。因此，宗徒在栽种、浇灌、建筑、立根基上工作，凡不愿这样工作的，就不应当吃饭。因为若不以天主之言语神粮来养育别人，自己从中得到什么益处呢？正如那懒惰的仆人，领受了一个塔冷通却埋藏起来，不为主人营利，最终不是被夺去，并被丢在外面的黑暗中吗？同样，他们说道，我们也这样做：我们与那些从世界忙碌中疲惫地来到我们这里的弟兄们一同诵读，使他们在我们这里，在天主的话语、祈祷、圣咏、赞美诗和属灵歌曲中得到安息。我们向他们讲话，劝慰、鼓励，在他们身上按他们的身份修造我们所察觉到的缺失。如果我们不做这些工作，我们将危险地从主那里领受属灵的食粮本身。这就是宗徒所说的：\BibleQuote{谁若不愿工作，就不应当吃饭}。这样，这些人以为自己既符合宗徒的教导，又符合福音的教导，因为他们认为福音吩咐我们不要顾虑此生的肉身和暂时的需要，而宗徒所说的\BibleQuote{谁若不愿工作，就不应当吃饭}是指属灵的工作和食物。\par

3. 他们也没有留意到：倘若另有人说，主确实以比喻和象征谈论属神的饮食和衣服，警告祂的仆人不应为这些事忧虑（正如祂说：\BibleQuote{当人把你们拖到审判台前时，不要思虑你们要说什么，因为在那一刻，必会赐给你们要说什么；其实说话的并不是你们，而是你们父的圣神在你们内说话。}因为祂不愿他们为属神智慧的话语忧虑，应许要赐给他们，使他们毫不挂虑）；但宗徒如今以宗徒的方式，更明确、更直接地（而非以比喻的方式，如同在他的宗徒书信中许多——几乎全部——那样）谈论体力劳动和食物时，明确地说：\BibleQuote{若有人不愿工作，就不该吃饭。}那些人会使自己的判断变得可疑，除非他们思想主的其他话语，找到证据证明祂说\BibleQuote{不要为吃什么、喝什么、穿什么忧虑}是指不为属肉体的食物和衣服忧虑。譬如，若他们留意祂所说\BibleQuote{这一切都是外邦人所寻求的}，因为祂在那里表明所指的是非常属肉体和暂时的事物。因此，倘若宗徒在此事上只说了这一句话：\BibleQuote{若有人不愿工作，就不该吃饭}，那么这句话或许可以被曲解为另一种意思；但由于他在书信的许多其他地方极为清楚地表明了他在这点上的心意，他们徒劳地试图在自己和别人面前制造迷雾，以便他们不仅拒绝实践那爱德所劝勉的，甚至拒绝自己理解或让别人理解它；他们不惧怕那记载的：\Quote{他不愿明白，好去行善。}\par

4. 首先，我们应当证明：蒙福的宗徒保禄愿意天主的仆人从事体力劳动，这些劳动应以伟大的属神赏报为目的，为使他们无需从任何人那里获取食物和衣服，而是用自己的双手为自己提供这些；然后，要表明那些福音的训诫（有些人不仅从中滋长懒惰，甚至滋生傲慢）与宗徒的训诫和榜样并不矛盾。现在，让我们看看宗徒是如何得出这样的结论，以至于他说：\BibleQuote{若有人不愿工作，就不该吃饭。}以及他随后补充的话，这样从这段经文的上下文中可以看出他明确的判决。\BibleQuote{弟兄们，我们奉主耶稣基督的名命令你们，要远离一切行为不端正、不按照我们所传给你们的传统而行的弟兄。因为你们自己知道应当怎样效法我们；我们在你们中间并不是不端正的，也没有白吃任何人的饭，而是辛苦劳碌，昼夜工作，免得加重你们任何人的负担：这并不是因为我们没有权柄，而是要给你们作榜样，叫你们效法我们。因为我们从前在你们那里时，曾吩咐过你们：若有人不愿工作，就不该吃饭。因为我们听说，你们中间有人行为不端正，什么工也不作，却好管闲事。我们靠主耶稣基督吩咐并劝勉这样的人，要安静地工作，吃自己的饭。}对这些话还能说什么呢？为了使此后无人能凭自己的意愿（而非按爱德）随意解释，他用自己的榜样教导了他所命令的。因为作为宗徒、福音的宣讲者、基督的士兵、葡萄园的栽种者、羊群的牧人，主已为他规定他应靠福音生活；但他自己并未索取应得的报酬，为给那些觊觎不应得之物的人立下榜样，正如他对格林多人所说：\BibleQuote{有谁当兵而自备粮饷呢？有谁栽种葡萄园，而不吃它的果实呢？有谁牧养羊群，而不喝羊奶呢？}因此，他拒绝接受应得之物，为借他的榜样制止那些在教会中未受如此任命却自以为应得同样待遇的人。他说：\BibleQuote{我们也没有白吃任何人的饭，而是辛苦劳碌，昼夜工作，免得加重你们任何人的负担；这并不是因为我们没有权柄，而是要给你们作榜样，叫你们效法我们。}这是什么意思？因此，那些他给予此命令的人——即那些没有他所具有的权柄、不能仅凭属神的工作而吃并非由体力劳动挣来的面包的人——应当听从他。正如他所说：\BibleQuote{我们靠主吩咐并劝勉，要安静地工作，吃自己的饭。}他们不应反驳宗徒最清楚的话，因为这也属于那\Quote{安静}，即他们应当安静地工作并吃自己的饭。\par

5. 然而，若不是我在他书信中发现了其他更为显明的经文，我本会以更为详尽和勤勉的态度来探讨和处理这些话语；藉着比较这些经文，上述两处经文便更为清晰地显明出来，而且即使没有这些经文，那些经文也足够。他向格林多人，也就是针对同一件事，这样写道：\BibleQuote{我不是自由的吗？我不是宗徒吗？我没有见过我们的主耶稣基督吗？你们不是我在主内的工作吗？如果我对别人不是宗徒，对你们我确实是宗徒，因为你们在主内正是我宗徒职务的印记。我的辩护是向那些质问我的人说的：难道我们没有权利吃喝吗？难道我们没有权利像其他的宗徒、主的弟兄和刻法一样，带领一位姊妹同行吗？}请注意，他首先表明什么是合法的，而且之所以合法是因为他是宗徒。因为他以此开头：\BibleQuote{我不是自由的吗？我不是宗徒吗？}并证明自己是宗徒，说：\BibleQuote{我没有见过我们的主耶稣基督吗？你们不是我在主内的工作吗？}这一点得到证明后，他便表明对他合法的事，对其他宗徒也是合法的；也就是说，他不应亲手做工，而应靠福音生活，正如主所规定的，他在下文最清楚地表明了这一点；因为为此目的，有属世财产的信女也跟随他们，并以她们的财产供给他们，使他们不缺少生活所需的一切。有福的保禄表明，这件事对他固然合法，正如其他宗徒所做的那样，但他后来提到，他没有选择行使这一权力。有些人不了解这一点，当他说\BibleQuote{难道我们没有权利带领一位姊妹女人同行吗？}时，他们解释为\InnerQuote{一位姊妹妻子}，而不是\InnerQuote{一位姊妹女人}。他们被希腊语一词的歧义所误导，因为希腊语中\OriginalTerm{妻子/女人}{\GreekText{γυνή}}一词既表示\Quote{妻子}，也表示\Quote{女人}。事实上，宗徒已经这样表达，使得他们不应犯此错误；因为他既没有只说\Quote{女人}，而是\Quote{姊妹女人}；也没有说\Quote{娶}（指婚姻），而是\Quote{带领}（指旅行）。然而，其他译者并未被这一歧义所误导，他们解释为\Quote{女人}，而不是\Quote{妻子}。\par

6. 若有人认为，圣洁的妇女陪伴宗徒到他们宣讲福音的任何地方，以她们的财产供给他们的需要，这是不可能的，就让他听听福音，学习他们在这一点上如何效法了主自己的榜样。也就是说，我们的主按祂仁慈的惯常做法，同情软弱的人，尽管天使可以服侍祂，祂却有一个钱袋，里面装着虔诚信友（无疑是为了生活所需）施舍的银钱，祂将这钱袋交给犹大保管，为的是让我们学会在教会中容忍窃贼，如果我们不能避免他们的话。因为正如经上论及他所说：\BibleQuote{他偷取其中所放的。}祂也愿意妇女跟随祂，为了预备和供应必需品，以此表明福音宣讲者和天主的仆人有如士兵，应从作为属民的天主子民那里得到当得的供给；这样，如果有人不愿使用当得之物，正如宗徒保禄不愿使用一样，他就能给予教会更多，因为他不索取当得的报酬，而是靠自己的劳动赚取日用的生活费。因为曾对那个带受伤者的人说过：\BibleQuote{你所额外花费的，我回来时必还给你。}因此，宗徒保禄确实是\Quote{额外花费了}，因为他如同他自己作证，自费作战。因为在福音中记载：\BibleQuote{此后，祂周游各城各乡，宣讲并传扬天主的国，十二门徒与祂同在，还有一些曾被恶灵和疾病治愈的妇女：号称玛利亚的玛达肋纳，从她身上曾赶出七个魔鬼；还有黑落德的家宰雇撒的妻子约安纳，以及苏撒纳，和许多别的妇女，她们都用自己的财产供给祂。}宗徒们效法了主的这个榜样，接受当得的食物；关于这一点，同一主最明确地说：\BibleQuote{你们去，宣讲说：天国临近了。医治病人，复活死人，洁净癞病人，驱逐魔鬼。你们白白得来的，也要白白施舍。不要在腰带里备有金银铜钱；路上不要带口袋，不要带两件内衣，也不要带鞋，也不要带棍杖；因为工人配得他的食物。}看，主正是吩咐了宗徒所提到的事。祂叫他们不带这一切，正是为了在需要时，他们可以从那些向他们宣讲天国的人那里得到供应。\par

7. 但恐怕有人以为这只是赐给了十二位宗徒，请看路加所记载的：\BibleQuote{这些事以后，主又拣选了其他七十二人，派遣他们两个两个地在祂前面，到祂将要去的各城各地。祂对他们说：庄稼固多，工人却少，所以你们应当求庄稼的主人，派遣工人来收割祂的庄稼。你们去吧！看，我派遣你们犹如羔羊往狼群中。你们不要带钱囊，不要带口袋，也不要带鞋，路上也不要向人请安。你们无论进了哪一家，先要说：愿这一家平安。那里如有和平之子，你们的平安就要停留在他身上；否则，你们的平安仍归于你们。你们要住在那一家，吃喝他们所供给的，因为工人自当得他的工资。} 由此可见，这些不是命令，而是许可，凡愿意使用的，就可以用主所准予的合法权利；但若有人不愿使用，他并非违反命令，而是放弃自己的权利，以更慈悲和劳苦的方式在福音中行事，甚至不接受他所应得的工资。否则，宗徒就违反了主的命令：因为他在表明那对他合法之后，立即补充说：\BibleQuote{然而我并没有用过这权利。}\par

8. 但让我们回到我们论述的顺序，并仔细思量书信的整段经文。他说：\BibleQuote{难道我们没有权利吃喝吗？难道我们没有权利娶一位姊妹为妻吗？} 他所说的权利是什么，岂不是主赐给那些被派遣宣讲天国的人的权利吗？祂说：\BibleQuote{你们吃他们所供给的，因为工人自当得他的工资；} 并把自己作为同样权利的榜样，最忠信的妇女们曾用自己的财物供应祂所需。但保禄宗徒做得更多，他从同宗徒们那里引证这主所许可的权利的证据。因为他不是指责而补充说：\BibleQuote{就如其余的宗徒、主的弟兄和刻法一样；} 而是借此表明，他之所以不接受，是因为那是他合法接受的，这可以由其他战友的习惯所证明。\BibleQuote{难道只有我和巴尔纳伯没有权利不做工吗？} 看，他消除了即使是最迟钝之人心中的一切疑惑，使他们明白他所说的是何种工作。因为他说：\BibleQuote{难道只有我和巴尔纳伯没有权利不做工吗？} 这是什么意思？岂不是所有福音传道者和天主圣言的仆人都从主那里获得了权利，不必亲手做工，而可以靠福音生活，只从事属灵的工作，即宣讲天国、建立教会的和平？因为没有人能说，保禄所说的\BibleQuote{难道我和巴尔纳伯没有权利不做工吗？} 是指那种属灵的工作。因为所有那些人都有不做工的权柄；那么，凡试图败坏和歪曲宗徒训导的人，让他说吧，如果他敢说所有福音传道者都从主那里获得了不宣讲福音的权利。但若这是极其荒谬和疯狂的说法，他们为何不明白显而易见的事呢？他们确实获得了不做工的权利，但不是指不做属灵的工作，而是指不做体力劳动以维持生计，因为\BibleQuote{工人自当得他的工资}，正如福音所说。因此，不是保禄和巴尔纳伯没有不做工的权利；而是所有的人都同等拥有这权利，但这些人没有使用这权利，\Quote{更多地付出}于教会，以便在他们宣讲福音的地方，他们认为对软弱的人适宜。为此，他为了不显得责备同宗徒们，继续说道：\BibleQuote{谁当兵而自备粮饷呢？谁栽种葡萄园，而不吃园中的果子呢？谁牧养羊群，而不喝羊群的奶呢？我说这些话，岂是按照人的意见呢？律法不是也这样说吗？因为在梅瑟的律法上记载：\OriginalTerm{牛在场上踹谷的时候，不可笼住它的嘴}{You shall not muzzle the ox that treads out the grain}。天主岂是关心牛吗？还是全为我们说的呢？的确，是为我们写的，因为犁地的应当怀着希望犁地，打场的也当怀着分享谷物的希望。} 保禄宗徒以这些话充分表明，同宗徒们并不曾越分，他们不做体力工作以获取生活必需品，而是照主的安排，靠福音生活，白白地从他们所宣讲白白恩宠的人那里接受免费的面包。他们如同士兵接受军饷，从他们所栽种的葡萄园的果实中，按需自由收取；从他们所牧养的羊群的奶中，他们喝；从他们所打场的禾场上，他们取食。\par

9. 但他接着下面所说的话更加明显，完全消除了所有疑惑的理由。\Quote{如果我们给你们，}他说，\Quote{播下了属灵的事物，我们若收取你们属血肉的事物，难道是大事吗？}他播下的属灵事物是什么，不就是天国之圣事的言语和奥迹吗？而他说他有权收取的属血肉事物，不就是那些为了肉身的生命和贫困而给予的现世事物吗？然而，他宣称这些虽是应得的，他并未寻求或接受，以免给基督的福音造成任何阻碍。那么，我们还能认为他做的是什么工作来维持生计呢？除了用他有形可见的双手做有形的工作之外，还有什么？因为如果他从属灵的工作中寻求饮食，即从那些他藉福音所建树的人那里接受这些，他就不能像他接下来所说的那样：\Quote{若别人在你们身上分享这权柄，我们岂不更可以吗？然而，我们没有用过这权柄，反倒凡事容忍，免得给基督的福音造成任何阻碍。}他说他没有用过的权柄是什么？不就是他从主所得的、对他们所拥有的权柄，即收取他们属血肉的事物，以维持这活在肉身中的生命吗？这权柄也是那些没有首先向他们传福音，而是后来来到他们的教会宣讲同一福音的人所分享的。因此，当他说了\Quote{我们若给你们播下了属灵的事物，我们若收取你们属血肉的事物，难道是大事吗？}之后，他接着又说：\Quote{若别人在你们身上分享这权柄，我们岂不更可以吗？}当他显示了他们所拥有的权柄之后，他说：\Quote{然而，我们没有用过这权柄，反倒凡事容忍，免得给基督的福音造成任何阻碍。}因此，让这些人说说，宗徒如何从属灵的工作中得到属血肉的食物，而他自己公开说他没有用过这权柄。但如果他从属灵的工作中得不到属血肉的食物，那么剩下的就是他从肉身的工作中得到，并且他对此说：\Quote{我们也没有白吃谁的饭，而是昼夜辛苦劳碌，免得你们任何人受累；并不是因为我们没有权柄，而是要给你们作榜样，叫你们效法我们。}他说：\Quote{我们凡事忍受，免得给基督的福音造成任何阻碍。}\par

10. 他又回过头来，一再反复强调他有权做却不做的事。\Quote{你们岂不知道，}他说，\Quote{在圣殿里供职的，就吃圣殿中的物；在祭台前服务的，就分享祭台上的物吗？主也这样吩咐，叫传福音的人靠福音养生。但这一切权柄我都没有用过。}还有什么比这更明显的？还有什么比这更清楚的？我怕我在谈论中想要解释这事，反而使本身明亮清楚的事变得晦暗。因为那些不明白这些话的人，或假装不明白的人，更不容易明白我的话，或假装明白；除非他们因此很快明白我们的，那是因为他们允许自己嘲笑明白的事物；但对于宗徒的话，这是不允许的。因此，凡他们不能按自己的意思另作解释的地方，尽管那是多么清楚明白，他们却回答说那是晦暗不明的，因为他们不敢说它是错误和乖戾的。天主的人大声说：\Quote{主也这样吩咐，叫传福音的人靠福音养生；但我一切权柄都没有用过。}血肉之躯却试图使直的变弯，开的变闭，清的变浊。它说：\Quote{他做的是属灵的工作，并以此为生。}如果是这样，他就是靠福音养生；那么他为什么说：\Quote{主也这样吩咐，叫传福音的人靠福音养生；但我一切权柄都没有用过？}或者，如果他们坚持也要把这\Quote{养生}一词解释为属灵的生命，那么宗徒对天主就没有希望了，因为他没有靠福音养生，因为他曾说：\Quote{我一切权柄都没有用过。}因此，为了使宗徒有永生的确定希望，他无论如何是属灵地靠福音养生。所以，他所说的\Quote{但我一切权柄都没有用过}无疑让我们理解到，关于主吩咐传福音的人靠福音养生的话，是指这肉身中的生命；即他们要靠福音维持这需要饮食和衣物的生命；正如上面他谈到同为使徒的人；主亲自说过：\Quote{工人配得他的食物；}以及\Quote{工人配得他的工钱。}因此，这食物和维持这生命的工钱，是应该付给传福音者的，宗徒却没有从他所传福音的人那里接受，他说的真话：\Quote{我一切权柄都没有用过。}\par

11. 他继续接着说，免得有人以为他之所以没有接受，只是因为未曾给予：\BibleQuote{但我写这些事，并非要人这样待我；我宁愿死，也不愿有人使我的光荣落空。}这是什么光荣？除非是他在基督里与软弱者同受苦时所愿在天主面前拥有的光荣。正如他即刻要极其坦白地说：\BibleQuote{我若传福音，并没有什么可夸的，因为责任已托付给我；}也就是说，维持此生的必要。\BibleQuote{我若不传福音，我就有祸了！}他说，意思是若我不传福音，我将因饥饿而受苦，没有生活所需。因为他继续说：\BibleQuote{我若甘心情愿做这事，便有赏报。}他所谓甘心情愿，是指他不受维持现世生命的任何必要所迫而去做；为此他得到赏报，就是在天主面前永恒的荣耀。\BibleQuote{但若不情愿，管家职务却委托给我了；}意思是，我若不情愿，却因度过现世生命的必要而被强迫传福音，\BibleQuote{管家职务委托给我了；}也就是说，由于我作为管家的职务，因为基督、因为真理是我所传的，无论动机如何，无论寻求自己的利益，无论因世俗利益的需要而被迫如此，别人得益处，但我却没有在天主面前那光荣而永恒的赏报。\BibleQuote{那么，我的赏报是什么？}他以提问的方式说；因此发音须停顿，直到他给出答案。为更容易理解，让我们仿佛向他提问：\BibleQuote{那么，宗徒啊，当你不接受那应得的、为传福音者所定的、并非为它而传却因主的规定而附带给予的世俗赏报时，你的赏报将是什么？你的赏报会是什么呢？}看他如何回答：\BibleQuote{那就是，传福音时，使基督的福音不费分文；}也就是说，福音对信徒不至于昂贵，免得他们以为福音传给他们是为了让传道者似乎是在出售它。然而他又一再回头，为要显明按主的规定他有权接受却没有接受：\BibleQuote{免得我滥用福音中的权柄。}\par

12. 但是，现在，他这样做是为了容忍人的软弱，让我们听听接下来的话：\BibleQuote{因为我虽是自由的，不属于任何人，却使自己成为众人的仆人，为要赢得更多的人。对在法律之下的人，我就成为在法律之下的，为要赢得那些在法律之下的人；对没有法律的人，我就成为没有法律的——其实我并不在天主面前没有法律，而是在基督的法律之下——为要赢得那些没有法律的人。}他这样做，并非出于诡诈的\OriginalTerm{伪装}{simulation}，而是出于\OriginalTerm{怜悯之同情}{compassion}；也就是说，并非如有些人所想的那样，他假装是\Person{犹太人}，因为他在\Place{耶路撒冷}遵守了旧法律所规定的事。因为他这样做是符合他自由而公开宣告的论断，他说：\BibleQuote{若有受割损而蒙召的，就不要成为未受割损的。}也就是说，不要让他生活得好像他成了未受割损的，并将他所暴露的遮盖起来；正如他在另一处所说：\BibleQuote{你的割损成了未受割损。}因此，正是符合他所说的这个论断：\BibleQuote{若有受割损而蒙召的，就不要成为未受割损的。有未受割损而蒙召的，就不要受割损；}他做了那些事，而那些事被那些不理解且不够注意的人视为他在\OriginalTerm{假装}{feign}。因为他本是\Person{犹太人}，且是在受割损时蒙召的；因此他不愿成为未受割损的；也就是说，他不愿生活得好像他未受割损。因为这事他现在有权去做。事实上，他并不像那些奴性地遵行法律的人那样在\OriginalTerm{法律之下}{under the law}；他是在天主和基督的\OriginalTerm{法律内}{in the law}。因为那法律并不是一种，而天主的法律是另一种，如同可咒骂的\OriginalTerm{摩尼派信徒}{Manicheans}惯常所说的那样。否则，如果当他做那些事时，他被视为在假装，那么他\OriginalTerm{假装}{feign}自己也是一个外教人，并向偶像献祭，因为他说他对那些没有法律的人，就成了没有法律的。毫无疑问，他指的无非就是我们称为外教人的外邦人。因此，在法律之下是一回事，在法律内是另一回事，没有法律又是另一回事。\OriginalTerm{在法律之下}{under the law}的是属血肉的犹太人；\OriginalTerm{在法律内}{in the law}的是属神的人，包括犹太人和基督徒（因此前者遵守了他们祖先的习俗，却没有把不寻常的重担加在相信的外邦人身上；而且他们也受了割损）；但\OriginalTerm{没有法律}{without law}的是尚未相信的外邦人，宗徒证明自己变得像他们一样，是因为同情之心，而不是外表变幻的伪装；也就是说，他可以用那种方式帮助属血肉的犹太人或外教人，正如他自己，如果他是那一种人，也会希望得到帮助一样：即承担他们的软弱，在相似的同理心之中，而不是在谎言的虚构中欺骗；正如他紧接着所说：\BibleQuote{对软弱的人，我就成为软弱的，为要赢得那软弱的人。}因为他正是在这一点上说话，说了所有那些其他事。因此，如同他成为软弱的人中的软弱者并非谎言，以上提到的所有其他事也是这样。因为他所说的对软弱者的软弱，若不是与他们一同受苦，又是什么呢？以至于，为了不让他看起来像是福音的贩卖者，并因陷入无知之人的恶意猜疑而阻碍天主圣言的道路，他拒绝接受主所应许给他当得之物。如果他愿意接受，他绝不会说谎，因为那确实是应得的；而因为他没有接受，他也绝没有说谎。因为他没有说它不该得；而是表明它该得，并且既然该得，他便没有使用它，并声称他根本不会使用它，就在这点上成为了软弱的；即他不使用自己的权力。他具有如此仁慈的情感，以至于他想到，如果他自己也变得如此软弱，他会希望别人如何对待他；可能，如果他看到那些向他宣讲福音的人接受他们的费用，他可能会认为这是贩卖商品，并因此猜疑他们。\par

13. 论及自己的软弱，他在另一处说：\BibleQuote{我们在你们中间变得渺小，有如乳母抚育自己的孩子。} 因为在那一处，上下文表明这一点：他说：\BibleQuote{因为我们从未用过谄媚之词，正如你们知道的，也未用过贪婪的借口；天主是见证。也没有寻求人的光荣，无论是从你们还是从别人，虽然作为基督的宗徒，我们可以加重你们的负担，但我们在你们中间变得渺小，有如乳母抚育自己的孩子。} 因此，他给格林多人所说的话——他具有宗徒的权力，如同其他宗徒一样，而他见证自己从未使用这权力——这同样也是他在那处给得撒洛尼人说的：\BibleQuote{虽然作为基督的宗徒，我们可以加重你们的负担。} 按照主所说：\BibleQuote{工人配得他的工资。} 因为这是他所谈论的事，由他上面所记的话表明：\BibleQuote{也未用过贪婪的借口，天主是见证。} 即是说，由于主所规定的权利本应归于善传福音的人，他们传福音并非为了这权利，而是寻求天主的国，这样这一切都加给他们了，然而有人却从中取利，关于这些人他也说：\BibleQuote{因为这样的人并非服事天主，而是服事自己的肚腹。} 宗徒希望断绝这些人的机会，以至于即使是合法应得的，他也放弃。因为这一点他在\BibleBook{格林多}{后书}中清楚表明，谈及别的教会供应他的需要。因为他似乎曾陷入如此极度的匮乏，以致远方的教会寄来供应他的所需，而在他所居留的人中，他却未接受任何这类东西。他说：\BibleQuote{我犯了罪吗？我贬抑自己，为使你们高举，因为我白白地向你们宣讲了天主的福音？我掠夺了别的教会，取了他们的薪俸，为服事你们。当我在你们那里，并有所缺乏时，我没有拖累过任何人。因为从\Place{马其顿}来的弟兄们补足了我的缺乏，在一切事上我保持自己不拖累你们，并且将继续保持。基督的真实在我内，这光荣在\Place{阿哈雅}地区不会在我身上受到限制。为什么？因为我不爱你们吗？天主知道。但我所做的，我也要去做，为断绝那些寻找机会之人的机会，好使他们所光荣的，也像我们一样，被找着。} 因此，关于这机会，他在这里说他断绝了，他愿意人理解他在前一处所说的：\BibleQuote{也未用过贪婪的借口，天主是见证。} 而他在这里所说的：\BibleQuote{我贬抑自己，为使你们高举}，这在给同一格林多人所写的第一封书中是：\BibleQuote{对软弱的人，我就成为软弱的}；在给得撒洛尼人书中是：\BibleQuote{我在你们中间变得渺小，有如乳母抚育自己的孩子。} 那么，现在请注意下文：他说：\BibleQuote{我们既然这样恋爱你们，就不仅愿意将天主的福音分给你们，连自己的灵魂也愿意分给你们，因为你们已成为我们极亲爱的人。弟兄们，你们记得我们的劳苦和艰辛，日夜工作，为的是不拖累你们任何人。} 因为他在上面说了：\BibleQuote{虽然作为基督的宗徒，我们可以加重你们的负担。} 因为那时软弱的人处于危险中，恐怕他们被虚假的猜疑所激动，憎恨一种似乎是可出卖的福音，为此原因，他以父母般的慈爱为他们战兢，做了这事。同样，在\BibleBook{}{宗徒大事录}中他也说了同样的话，当他从\Place{米肋托}派人到\Place{厄弗所}，从那里请来教会的长老们，对他们，除了许多别的话之外，他说：\BibleQuote{我没有贪图过任何人的金银或衣物；你们自己知道，我的双手供应了我自己的需要和同我在一起的人的需要。在一切事上我都给你们指出，这样劳苦应当扶助弱小，并记住主耶稣的话，因为他说：\InnerQuote{施予比领受更为有福。}}\par

14. 也许有人会说：\Quote{如果宗徒从事体力劳动以维持生计，那么他做的是什么工作，又如何在传福音的同时找到时间做工呢？}我回答：即使我不知道，但以上所述无疑都证明他确实从事体力劳动，以此在世度日，并且没有使用主赐给宗徒的权利，即传福音的人应靠福音生活。因为这些话并非只在某处或简略提及，以至于任何最狡猾的诡辩者都无法用尽一切遁词来曲解和歪曲其含义。既然如此伟大的权威以如此有力而频繁的打击击溃反对者，击碎他们的悖逆，他们又何必问我他做了何种工作，或何时做工？我只知道，他没有偷窃，没有做盗贼、强盗、车夫、猎人、戏子或贪图不义之财，而是清白正直地从事有益于人类使用的工作，如木匠、建筑工、鞋匠、农夫等类似的工作。因为正直本身并不责备那些骄傲之人所责备的事，这些人喜欢被称为正直，却不喜欢成为正直。那么，宗徒不会不屑于从事任何农夫的工作，也不会不屑于工匠的劳作。因为那说\BibleQuote{你们不可使犹太人、希腊人或天主的教会跌倒}的人，我不知道他在什么人面前可能感到羞愧。如果说是犹太人：圣祖们曾牧放牲畜；如果说是希腊人，我们也称之为外教人：他们有备受尊崇的哲学家是鞋匠；如果说是天主的教会：那位义人，被选为永恒童贞的见证，与生育基督的童贞玛利亚订婚的，正是一个木匠。因此，凡清白无欺的工作都是善的。宗徒自己也提防这一点，使人不因维持生计而转向恶行。他说：\BibleQuote{那偷窃的，不要再偷；却要劳苦，亲手做善事，好能周济有需要的人。}因此，知道宗徒在体力劳动中从事善事就已足够。\par

15. 但他常在什么时候工作，即在什么时间做工而不妨碍传福音，谁能确定呢？不过，他确实日夜工作，他自己也未曾隐瞒。然而，这些人仿佛十分忙碌而问及工作时间，他们做了什么？他们是否从\Place{耶路撒冷}直到\Place{依里黎果}，用福音充满了大地？或者，还有哪些蛮族地区未被他们前往，以教会和平来充满？我们知道他们聚集成某种神圣的团体，极其悠闲。宗徒做了一件奇妙的事：在他对所有教会（无论是已建立的还是待建立的）如此巨大的关怀中，他仍然亲手做工；但正因如此，当他与\Place{格林多}人在一起时，尽管有缺乏，却没有成为任何当地人的负担，而来自\Place{马其顿}的弟兄们补充了他所欠缺的一切。\par

16. 他本人也顾及圣徒们类似的需用，他们虽然遵守他的诫命，\BibleQuote{安静作工，吃自己的饭}，但仍可能因诸多原因在类似生计方面有所欠缺。因此，在他如此教导和劝诫之后说：\BibleQuote{现今我们因主耶稣基督命令并劝勉这样的人，要安静作工，吃自己的饭；}然而，惟恐那些有能力供应天主仆人之需的人因此变得懒惰，他随即加上：\BibleQuote{至于你们，弟兄们，行善不可懈怠。}他在写给弟铎的信中说：\BibleQuote{你要尽力送则纳律师和阿颇罗前行，使他们一无所缺；}为表明他们应在何事上一无所缺，他立刻续道：\BibleQuote{但我们也该学会行善，以应付急需，免得不结果实。}至于弟茂德，他称他为最真实的儿子，因他知道他身体软弱（如他在劝他不要只喝水，为胃和屡次患病可稍用酒时所显示的），惟恐他因身体不能作工，又不愿从他所传福音的人手中接受日用的食物，而寻求某种使心神陷入纠缠的事业（因为心神自由而身体劳作是一回事，如一个手艺匠，若他不欺诈、不贪婪、不贪图私利；但另一回事，是操劳心神以聚集钱财而无身体劳作，如商人、管家或承办者，他们用心思的关切处理事务，而非亲手作工，并且在这方面以谋利的焦虑占据心神本身）；惟恐弟茂德因身体软弱不能亲手作工而堕入此种行径，他这样劝勉、告诫并安慰他：\BibleQuote{你要努力，如同耶稣基督的精兵。没有人当兵为天主而战，会纠缠于世俗的事务，为要讨那招选他的人喜悦。凡比武的，除非按规矩竞赛，得不到花冠。}因此，惟恐那人口中说：\BibleQuote{我不能挖，乞讨又羞耻，}他又续道：\BibleQuote{劳力的农夫应先尝果实。}正如他对格林多人所说的：\BibleQuote{谁曾自备粮饷去当兵呢？谁栽种葡萄园而不吃园中的果实呢？谁牧养羊群而不喝羊的奶呢？}这样，他使一位纯洁的福传者无忧无虑；他作福传者并非为出卖福音，但也没有力量亲手供给自己生活所需；因为他应当明白，凡他所需要的，从他作为士兵所服务的省民、作为葡萄园所栽培、作为羊群所牧养的人那里收取的，不是乞丐的所得，而是权能的所用。\par

17. 因此，或因天主仆人的这些工作，或因身体上的软弱（这些无法完全避免），宗徒不仅允许善信信徒供应圣徒的需要，而且极其有益地加以劝勉。因为，他放下自己说的未曾使用过的权柄（但忠信者仍须服从这权柄），命令说：\BibleQuote{凡在圣言上受教的人，应把一切美物与施教的人分享。}放下这权柄（即圣言的宣讲者对所宣讲的人拥有的权柄），他常作见证；此外，他谈到那些变卖所有、分施一切、在耶路撒冷过着神圣共融生活的圣徒，他们不说任何东西是自己的，一切公用，心灵和灵魂在主内合一：他嘱咐并劝勉外邦人的教会应将他们所需之物供给这些人。因此，致 \BibleBook{}{罗马书} 中也说：\BibleQuote{所以，我现在往耶路撒冷去，为圣徒服务。因为马其顿和阿哈雅乐意捐出一笔款项，为供给耶路撒冷圣徒中的穷人。他们乐意这样作，而且他们也是他们的债务人。因为外邦人既然在他们属神的事上有了分，也就该在属肉体的事上供养他们。}这类似于他对 \BibleBook{格林多}{前书} 所说的：\BibleQuote{如果我们给你们播种了属神的事，若我们收割你们属肉体的事，还算大事吗？}同样，在 \BibleBook{格林多}{后书} 中：\BibleQuote{此外，弟兄们！我们将天主赐给马其顿教会的恩宠告诉你们：他们在患难大考验中的喜乐和极度的贫穷，竟涌出慷慨施舍的丰富。我为他们作证：他们是自愿的，按他们的力量，甚至超过力量，再三恳求我们，准他们分享供应圣徒的恩典和共融。他们所做的，不仅如我们所盼望的，而首先他们把自己献给主，也因天主的旨意献给我们。因此，我们请提托斯，既已开始，也要在你们中完成这恩宠的工作。但是，你们既然在一切事上，在信德、言论、知识、各样的热忱和你们对我们的爱德上，都丰盈有余，就要在这恩宠上也富足起来。我说这话不是命令，而是藉着别人的热忱，试验你们爱德的真诚。因为你们知道我们的主耶稣基督的恩宠：祂本是富有的，为了你们却成了贫困的，好使你们因祂的贫困而成为富有的。关于这事，我给你们建议：这为你们是有益的，你们从去年不但已开始实行，而且也有此心愿。如今就当完成这事，使你们既有心愿，也有实行，照你们所有的去完成。因为如果有甘心情愿，蒙悦纳是按人所拥有的，不是按他所没有的。这并非要使别人轻松，你们困难，而是出于均等：在这时刻，你们的丰盈补足他们的缺乏，好使他们的丰盈也成为你们缺乏的补足，以达到均等，正如经上所载：\InnerQuote{多收的没有余，少收的也没有缺。}但感谢天主，把同样的热忱为你们放在提托斯心里：他接受了劝勉，但更加热心，自愿前往你们那里。我们打发一位弟兄与他同去，这弟兄在众教会中因福音而受称赞；不仅如此，他也被众教会选派，作我们旅途的同伴，为这项由我们经手的恩宠工作，以光荣主，并显明我们的热忱：避免有人在这项由我们经手的丰盈事工上指责我们。因为我们关心善事，不但在主面前，也在人面前。}在这些话中，可以看出宗徒多么愿意圣善的团体不仅关怀供应圣善的天主仆人的必需品（他在这事上给予建议，因为这更有益于施予者本人，而非受惠者——因为对后者，另一样事有益，即他们应圣善地使用弟兄们对他们的这种服务，而非为此目的服侍天主，并取用这些物品仅为供给必需，而非滋养懒惰），而且蒙福的宗徒也说他亲自多么关注这项通过提托斯传递的事工，以致他告诉我们，一位同行的旅伴因此被众教会选派——一位天主的人，有很好的声誉，\BibleQuote{他的名声在众教会中因福音而传扬。}他接着说，同一人被选派为他的同伴，是为了避免人的指责，以免没有与他一起从事这事工的圣徒作证，软弱的和不信的人会以为他为自己收取并中饱私囊，而他所收取的却是为了供应圣徒的必需品，由他带来并分施给穷人。\par

18. 稍后他又说：\Quote{关于供应圣人的事，我写信给你们原是多余的，因为我知道你们的热心，我曾为此向马其顿人夸耀你们，说阿哈雅从去年就预备好了；而你们的热心也激动了很多人。但我还是派了那几位弟兄去，免得我们在这事上对你们的夸耀落空，好照我所说的预备妥当；万一马其顿人和我一同来到，见你们没有预备，我们（且不说你们）在这件事上就会羞愧。因此，我认为必须劝那几位弟兄先到你们那里，把你们从前所应许的捐款预备好，使这捐款成为慷慨的捐助，而不是出于吝啬。但我要说：少种的少收，多种的多收。各人要照心里所决定的捐助，不要为难，不要勉强，因为天主喜爱乐意捐献的人。天主能丰厚地赐予你们各种恩宠，使你们在一切事上常十分充足，能多多行各种善事，正如经上记载：\InnerQuote{他施舍钱财，周济穷人；他的仁义永存不废。}那赐种给播种者、赐粮给人作食物的，必供给并增多你们的种子，增加你们仁义的果实，使你们凡事富足，以致十分慷慨，这慷慨藉着我们使人感谢天主；因为这供应的事不但补足圣人的缺乏，而且使许多人藉着对天主的感谢而越发丰满，他们因着这供应的事的明证，就为你们承认基督福音的服从，以及你们对他们和众人慷慨的捐献，而光荣天主；他们也因天主在你们身上的极大恩宠而切切想念你们，为你们祈祷。感谢天主，为他不可言喻的恩赐！}宗徒如此谈论基督的士兵与其他臣民之间互相供给需要——一方面是以属肉体的事供给那些人，另一方面是以属神的事供给这些人——他是在何等丰盛的神圣喜乐中沉浸，以致像他那样呼喊，并仿佛因圣善喜乐的满溢而迸发出\Quote{感谢天主，为他不可言喻的恩赐！}\par

19. 因此，宗徒——不，更确切地说，那占据、充满并推动他心灵的圣神——不停地劝勉那些拥有这等财产的忠实信友，不要缺少那些愿意在教会中持守更高圣德境界的人的需要；这些人斩断了一切世俗希望的羁绊，并将自由的心志奉献于他们虔诚的作战服务。同样，这些人自己也应当服从宗徒的训令，同情软弱的人，不被私有财产的爱所束缚，亲手劳作以谋求公共福祉，并毫无怨言地服从他们的长上；这样，从善信者的奉献中为他们补足那因某些人的身体软弱、因教会职务或那带来救恩的教训的修养而仍然有所欠缺的部分——尽管他们劳作并从事某种工作以维持生计，但仍有不足。\par

20. 因为这些人——那些不愿做体力劳动的人——把时间花在什么事情上，我很想知道。\Quote{用在祈祷、咏唱圣咏、阅读和天主圣言上。}他们如此说。这无疑是圣善的生活，在基督的甘饴中值得赞美；然而，如果我们不能从中被召叫出来，那么我们就不能进食，也不能准备每日所需的食物，摆在我们面前并取用。既然天主的仆人在某些时刻因身体软弱被迫必须抽时间做这些事，为什么我们不为遵守宗徒的训令也分配一些时间呢？因为一个服从之人的一次祈祷比一万个轻慢之人的祈祷更易蒙垂听。至于神圣的歌曲，即便在用手劳作时，也容易吟唱，如同船夫唱着船歌，以神圣的旋律鼓舞他们的劳动。难道我们不知道所有工匠的情况吗？当他们手不停工时，有多少人会将自己的心思和言语投入到戏院故事的虚妄甚至污秽中去？那么，有什么能阻止天主的仆人一边用手劳作，一边默思上主的法律，歌颂至高上主的名呢？当然，为了能记住所学的内容并加以复述，他需要有专门的时间。为此，信友们那些善行也不应欠缺，以补足必要的生活所需，使那些用于充实心智而无法进行体力劳动的时辰不至于陷入匮乏。然而，那些声称把时间用于阅读的人，难道他们在阅读中找不到宗徒所命令的事吗？那么，这是何等的不合理：因想花时间阅读，就拒绝受他所读之物的管束；而为了能花更多时间阅读有益之事，便拒绝去做所读之事？因为谁不知道，一个人阅读善事时越快地吸收，他在实践所读之事时就越迅速呢？\par

21. 再者，倘若必须对某人施以讲论，且这讲论如此占据说话者，以致他没有时间亲手劳作，那么隐修院中所有弟兄都能对从另一种生活方式来到他们那里的弟兄们进行讲论，无论是讲解圣经课程，还是针对可能提出的任何问题以有益的方式进行推理吗？既然并非人人都有此能力，为何以此为借口，人人都想无所事事呢？即使人人都有能力，他们也应轮流进行；这不仅是为了使其余的人不被从必要的工作中抽离，也是因为对许多听众而言，有一位说话者就足够了。现在谈到宗徒：他如何能抽出时间亲手劳作，除非他为传布天主的圣言设定了某些固定时间？的确，天主也不愿这事对我们隐藏。因为宗徒是何种工匠，以及他在哪些时间忙于分施福音，圣经并未略而不提。即当他因启程之日迫近而匆忙时，在特洛阿，恰在一周的第一天，弟兄们聚集一起擘饼之际，他是如此热切，并且争辩如此必要，以至他的讲论延长到半夜，仿佛他们忘记了那天并不是斋日；但当他在某处停留较久并每日争辩时，谁能怀疑他为这项职务设定了某些固定的时辰呢？因为在雅典，由于他在那里发现了最勤于探求事物的人，关于他这样记载说：\BibleQuote{因此他每日在会堂里同犹太人辩论，并在市场上同那里的外邦居民辩论。}并非每日在会堂里，因为他在那里惯常在安息日讲论；但他说：\BibleQuote{在市场上，}他说，\BibleQuote{每日；}无疑是由于雅典人的好学。接着如此记载：\BibleQuote{有些伊壁鸠鲁派和斯多葛派的哲学家与他讨论。}不久之后，又说：\BibleQuote{雅典人和侨居那里的客人，除了谈论或听取一些新鲜事物之外，没有别的事可做。}假设他在雅典的那些日子都没有工作：为此，他的需要确实是从马其顿供应的，正如他在\BibleBook{格林多}{后书}中所说的：然而事实上，他既可在其他时辰也可在夜间工作，因为他身心皆如此强健。但当他离开雅典后，让我们看看圣经说什么：\BibleQuote{他每逢安息日在会堂里辩论；}这是在格林多。然而在特洛阿，由于他离开迫近，讲论延长至半夜，那是一周的第一天，称为主日：由此我们明白他不是与犹太人而是与基督徒在一起；那时叙述者本人也说他们聚在一起擘饼。这确实是最好的管理，即一切事物按时间分配并有秩序地进行，以免陷入令人困惑的纠缠之中，使我们的心神混乱。\par

22. 同样也记载了宗徒从事何种工作。\BibleQuote{此后，他离开雅典，来到格林多；找到一个名叫阿桂拉的犹太人，是本都人，新近从意大利来，并有他的妻子普黎史拉；因为克劳狄命令所有犹太人都离开罗马，他就投奔他们；因他与他们同业，便住在他们那里，一同做工：原来他们是制帐棚的。}如果他们试图对此作寓意解释，便显示他们在教会学识上是何等精通——他们以此夸耀自己把所有时间都花在这上面。至少，关于上面所引的那些话：\BibleQuote{只有我和巴尔纳伯，难道没有权柄不做工吗？}以及\BibleQuote{我们没有用过这权柄；}以及\BibleQuote{虽我们作为基督的宗徒，本可叫人受累；}以及\BibleQuote{昼夜做工，免得累着你们任何人；}以及\BibleQuote{主也这样吩咐，传福音的人应靠福音生活；但我没有用过这些权利：}以及这类其余的话，要么让他们另作解释，要么如果被真理最清晰的光辉所折服，让他们明白并服从；或者如果他们不愿或不能服从，至少让他们承认那些愿意的人更好，那些也能的人比他们更有福。因为，为自己身体的软弱辩护——无论是真实声称还是虚假假托——是一回事；但如此受欺骗又如此欺骗人，以至懒惰竟能在一群无知的人中间获得统治权，这甚至被认为是天主的仆人更大力地获得义德的明证，那是另一回事。因为，一个表明真实身体软弱的人，必须以人情对待；一个假托虚假软弱而无法被定罪的人，必须交给天主：然而，他们中没有一个树立了有害的榜样；因为一位善的天主仆人，既服侍他明显软弱的弟兄；当另一个欺骗时，如果他因不认为对方是坏人而相信他，他并不会为了成为坏人而效仿他；如果他不相信，他视之为欺骗者，但也不效仿他。但当一个人说：\Quote{这就是真正的义德：不做任何体力工作，效法空中的飞鸟，因为谁做任何这样的工作，就是违背福音。}那么，凡心智软弱的人听见并相信这话，那人——不是因为他如此度过所有时间，而是因为他如此犯错——必须为之哀哭。\par

23. 由此产生另一个问题；或许有人会说：那么，其他宗徒、主的弟兄和\Person{刻法}，因为他们没有做工，难道犯罪了吗？或者他们给福音造成了阻碍，因为蒙福的\Person{保禄}说他故意没有使用这权柄，免得给基督的福音造成任何阻碍？因为如果他们因不做工而犯罪，那么他们就没有领受过不做工、反而靠福音生活的权柄。但如果他们领受了这权柄，按主的命令，那传福音的人应当靠福音生活；并因祂的话：\BibleQuote{工人配得他的食物}；\Person{保禄}稍加伸展，不愿使用这权柄；那么他们确实没有犯罪。如果他们没有犯罪，他们就没有造成阻碍。因为造成福音的阻碍不被认为不是罪。如果这样，他们说：\Quote{对我们来说，}他们说，\Quote{也可以自由地使用或不使用这权柄。}\par

24. 我应当简要解决这个问题，如果我说——因为我也应当公正地说——我们必须相信宗徒。因为他本人知道，为何在外邦人的教会中，不宜传扬可买卖的福音；这不是指责他的同僚宗徒，而是区分他自己的职务；因为他们无疑在圣神的劝告下，这样分配了传福音的区域，使\Person{保禄}和\Person{巴尔纳伯}到外邦人那里，而他们到受割礼者那里。但他将这诫命赐给了没有同样权柄的人，这已经由前述诸多事显明了。然而，就我判断，我们的这些弟兄们轻率地将这种权柄归给自己。因为如果他们真是福音传播者，我承认，他们拥有它；如果他们是祭台的仆人、圣事的分施者，那么这当然不是僭取，而是对权利的明确维护。\par

25. 如果他们至少曾在这个世界拥有可轻易不靠手工劳动维持此生所需的财富，当这财富在他们皈依天主时分施给了穷人，那么我们就必须相信他们的软弱，并予以容忍。因为通常这类人，不像许多人认为的受过较好教养，而是事实上（这才是真相）受过较为柔弱教养，无法承受体力劳动的辛劳。在\Place{耶路撒冷}或许就有许多人如此。因为经上记载，他们变卖了房屋田地，将价银放在宗徒脚前，以便按各人所需分施。因为他们被发现就近处，且对外邦人有益，这些从远处、从偶像崇拜中被召来的人，正如经上所说：\BibleQuote{法律将出自熙雍，上主的话将出自耶路撒冷}，因此宗徒称外邦的基督徒为他们的债务人：\BibleQuote{他们是他们的债务人}，他说，\BibleQuote{因为，如果外邦人在他们的属灵事上已有分，他们也该在属肉身的事上供应他们。} 但现在进入这服务天主的圣召中的，有来自奴隶身份的、或自由人的、或由此被主人释放或将要释放的人，同样来自农夫的生活，以及来自工匠的职业与平民劳动的人，他们的教养无疑越艰苦对他们越有益：接纳他们不是大罪。因为许多这一类人后来成了真正伟大且值得效法的人物。因这缘故，\BibleQuote{天主拣选了世上的软弱，以羞辱那强壮的；又拣选了世上的愚拙，以羞辱那聪明的；并拣选了世上卑贱的，以及那无有的，如同有的，为使那有的归于无有，好叫无人在天主面前夸口。} 这种敬虔圣洁的思想，因此也使得那些没有带来生活改变证据的人也被接纳。因为我们不清楚他们到底是为了服务天主而来，还是因逃避贫穷劳苦的生活而空腹前来，想得食得衣；不仅如此，还要受那些他们一向被鄙视践踏之人的尊荣。这类人既然不能以身体软弱为借口而不做工，因为他们过去的生活习惯证明他们并非如此，便躲藏在一面不良学问的幌子下，企图从误解的福音中曲解诫命：真是\BibleQuote{天空的飞鸟}，却因骄傲而高升；又是\BibleQuote{田野的野草}，却体贴肉体。\par

26. 正是如此，宗徒说在那些不守规矩的年轻寡妇身上所发生的事，是必须避免的：\BibleQuote{同时她们又习于懒惰，串门踏户，不仅懒惰，而且饶舌，多管闲事，说些不该说的话。}这正是他对恶毒妇女所说的话，我们也在恶毒男人身上哀叹悲悼；这些人反对他——那人在他的书信中我们读到这些事——却因懒惰而饶舌，说些不该说的话。如果其中有任何人当初是以这样的目的进入圣战，为讨他所献身的那位的喜悦，那么，当他们身强体健、健康无恙时，不仅能够受教，而且按照宗徒的教导也能工作；但他们接受了这些懒惰而腐败的言论（由于他们初学无知，无法判断），就因有毒的感染而变得同样有害：不仅不效法那些安静工作的圣人们的服从，也不效法其他隐修院在健全纪律中按宗徒规则生活的方式，反而侮辱比他们更好的人，宣扬懒惰是福音的守护者，指责怜悯是违背福音者。因为对于软弱者灵魂的慈悲工作，远比向饥饿者分饼以照顾其身体更为慈悲。因此，但愿天主使这些想闲着手的人也完全闲着舌头。因为他们如果只是懒惰的例子，而且也是沉默的，就不会有那么多愿意效法他们的人了。\par

27. 然而，实际上，他们竟反对基督的宗徒，朗诵基督的福音。因为懒惰者的工作如此奇妙，他们竟想以福音本身来证明懒惰有理，而宗徒正是为了不让福音受到阻碍才命令并实践工作的。然而，如果我们从福音本身的言辞迫使他们按他们理解的方式生活，他们首先会设法说服我们，他们并不像他们所理解的那样去理解。因为他们当然说，他们不应该工作，因为天上的飞鸟既不播种也不收割，主给了我们这样的比喻，叫我们不要为这些必需品挂虑。那么他们为什么不注意后面的话呢？因为不仅说\BibleQuote{它们不播种，也不收割}，还加了\BibleQuote{也不收藏在\OriginalTerm{仓房}{apothecæ}里}。而\Quote{仓}可以称作\Quote{谷仓}，或者字面上说\Quote{仓库}。那么这些人为什么想要闲着双手而拥有满满的仓库呢？为什么他们储存并保留从别人劳动中得来的东西，以便每天都有所供应？简而言之，为什么他们碾磨和烹饪？飞鸟并不做这些。或者，如果他们找到一些人，他们可以说服这些人也做这项工作，即每天将做好的食物带给他们；但至少他们的水，他们要么从泉源取来，要么从蓄水池和井中打来并储存；飞鸟并不做这些。但是如果他们喜欢这样，那就让善良的信徒和最虔诚的永恒之王的臣民们努力将他们的服务延伸到国王最勇敢的士兵身上，甚至到他们不必为自己打满一壶水的程度——如果今天的人们已经在一种新的正义等级上超越了当时在\Place{耶路撒冷}的人，超越了他们的步伐。因为当时由于饥荒即将来临，并由当时的先知预言，善良的信徒从\Place{希腊}送来了粮食供应；我猜想他们用这些粮食做了面包，或至少让人做了；飞鸟并不做这些。但是，如果今天这些人，如我开始所说，在某种正义等级上已经超越了这些人，并在维持此世生活方面完全像飞鸟一样；那就让他们向我们展示人们向飞鸟所做的服务，如同他们希望别人为他们所做的那样——除非是那些被捕获并关在笼子里的飞鸟，因为它们不被信任，以免它们飞走就不回来；然而它们宁愿享受自由并从田野里得到足够的食物，也不愿吃人放在它们面前并准备好的食物。\par

28. 在这里，这些人又将在一个更崇高的正义等级上被超越，那些将如此安排自己的人：他们每天都要到田野里去，仿佛去吃草，找到什么就吃一顿，平息饥饿后返回。但显然，由于田野的看守人，如果主愿意赏赐翅膀就好了，这样天主的仆人即使被发现在他人的田里，也不会被当作小偷抓走，而是像椋鸟一样被赶走。然而，这样的人会尽力模仿飞鸟，而捕鸟人却无法抓住它。但是，看，让所有人允许天主的仆人这样：他们愿意时，就进入他们的田地，然后毫无畏惧地、精神焕发地离开；正如法律给\Place{以色列}人民规定的：不可抓住在田里的小偷，除非他想从那里带走任何东西；如果他只拿了他所吃的，他们就让他自由地、不受惩罚地离开。因此，当主的门徒掐麦穗时，犹太人是在安息日的问题上而非偷窃上指责他们。但是在一年中的那些时节，当田里找不到可以当场取食的食物时，该怎么办呢？任何人若试图将任何东西带回家，以便可以通过烹饪为自己准备食物，按照这些人的理解，他就会从福音中被指责说：\BibleQuote{放下它，因为飞鸟并不这样做。}\par

29. 但我们也承认这一点：一年到头，田地里总能找到树木、蔬菜或各种根茎，可以生吃；或者，至少可以如此锻炼身体，以至于需要烹饪的食物即使生吃也无害，人们甚至在严寒的冬天也能出去寻找食物；这样，就没有什么东西需要取来准备，也没有什么东西为明天积存。然而，那些数日不见人影、不让任何人接近、将自己关闭起来、虔诚祈祷生活的人，却不能遵守这些规则。因为他们习惯将自己关闭起来，同时储存一些最容易获得、最便宜的食品，足以应付他们打算不见任何人的那些日子；飞鸟并不这样做。至于这些人在如此奇妙的节制中操练自己，因为他们有闲暇做这些事，并且不是出于骄傲的自我表现，而是出于慈悲的圣善，将自己呈现给人效法，我不仅不责备，反而不知道该如何充分赞美他们。然而，按照这些人对福音言语的理解，我们该如何评价这样的人呢？或者，他们越圣洁，就越不像飞鸟？因为除非他们为自己储备多日的食物，否则他们无法像那样关闭自己？然而，对他们也对我们说：\BibleQuote{所以不要为明天忧虑。}\par

30. 因此，为了简括整个问题，让这些以误解福音来图谋曲解诫命的人，要么像天空的飞鸟一样不为明天忧虑；要么听从宗徒，如同亲爱的子女：是的，他们应当两者都做，因为两者并不矛盾。因为耶稣基督的仆人保禄绝不会劝人行与他主相反的事。那么，我们公开向这些人说：如果你们如此理解福音中的天空飞鸟，以至于不愿亲手劳作以获取食物和衣服；那么你们也不应为明天积存任何东西，就像天空的飞鸟一样不积存任何东西。但如果为明天积存一点东西，可能并不违反福音，因为经上说：\BibleQuote{你们看天空的飞鸟，它们不播种，也不收割，也不积蓄在仓里}；那么，通过体力劳动来维持肉身的生命，也可能并不违反福音，也不违反天空飞鸟的比喻。\par

31. 因为如果他们被福音催促，不应为明天积存任何东西，他们最正确地回答说：为什么主自己有一个口袋来积存收集的钱？为什么在即将发生饥荒时，提前那么久就将粮食供应送给圣祖们？为什么宗徒们如此为圣徒的贫困提供所需，以免以后缺乏，以至最有福的保禄在致格林多人的书信中这样写道：\BibleQuote{关于为圣徒捐款的事，我怎样吩咐了迦拉达众教会，你们也当照样做。每週的第一天，你们每人要照自己的进项抽出来留着，免得我来的时候才现凑。及至我来到，你们写信举荐谁，我就打发他们，把你们的捐资送到耶路撒冷去。若我也该去，他们可以与我同去。} 他们极其丰富且真实地引用了这些和其他许多例子。我们回答他们说：你们看，虽然主说：\BibleQuote{不要为明天忧虑}，你们并没有被这些话约束而不为明天保留任何东西；那么为什么你们说，因为同样的话，你们被约束什么也不做？为什么天空的飞鸟不为你们保留任何东西的榜样，而你们却要它们作为什么也不做的榜样？\par

32. 有人会说：\Quote{那么，天主的仆人，如果离开他从前在世俗中的行业，皈依灵修的生活和战争，却仍然必须像一个普通工人那样从事劳作，这对他有什么益处呢？} 就好像真的能轻易用言语说明，主在回答那位寻求永生计策的富人时，告诉他若愿成全当如何行，即变卖所有、分给穷人、跟随祂，这有多大的益处。或者，有谁像那位说\Quote{我奔跑不是徒然，劳苦也不是枉然}的人一样，以如此无阻的步履跟随了主？他既命令这些工作，也亲自实行。我们受如此伟大的权威教导和塑造，这应当足以作为我们舍弃旧有资财、亲手劳作的榜样。但我们也蒙主亲自助佑，或许能在某种程度上领会，天主的仆人离开他们从前的行业，却仍如此劳作，这究竟有什么益处。因为如果一个人从富有皈依这种生活方式，且身体并无软弱，我们难道如此缺乏基督的馨香，以致不明白：当他削除昔日使灵魂致命膨胀的过剩之物，为了获得今世生活所必需的那一点，甚至不拒绝一个普通工人卑微的劳作时，这是何等医治旧日骄傲的良药？如果一个人是从贫穷的境况皈依这种生活方式，那么如果他不再追求增加自己微薄私产的欲望，如今不再求自己的利益，只求耶稣基督的事，将自己转入共同生活的爱德，与那些在天主面前同心合意的人共处，以致无人说任何东西是自己的，而大家共用一切，他就不可认为自己仍在做从前的事。因为如果在这尘世的共和国中，其首领们（正如他们自己的文人常以最热烈的言辞所讲述的那样）如此将整个城市和国家全体人民的公共利益置于自己的私事之上，以致其中一人为征服\Place{阿非利加}而获得凯旋，却毫无所得为女儿置办嫁妆，除非元老院下令从国库中拨给嫁妆；那么，对于身为永恒之城、天上\Place{耶路撒冷}之公民的人，他应当如何对待自己的共和国呢？岂不应将他亲手劳力所得与弟兄共享，若弟兄缺乏什么，就从公共储备中供给，效法那位他遵循其诫命和榜样的人所说：\Quote{似乎一无所有，却拥有一切。}\par

33. 因此，那些舍弃或分施了从前富足或多少丰裕的资产，以虔诚而有益的谦卑选择被列于基督的穷人之中的人，如果他们身体强健且无教会职务（尽管他们带来了如此伟大的意向证明，并将从前所有的颇多或不少财物用于同一团体的贫困者，公共基金本身和弟兄之爱理应回报他们生活的供养），但他们若也亲手劳作，以除去来自更卑微生活境况、因此更习惯于劳作的懒惰弟兄的一切借口，那么他们这样做，比之将他们所有的财物分给穷人更为慈悲。如果他们不愿这样做，谁又敢强迫他们呢？那么，应当为他们在隐修院内找到某些工作，这些工作虽较少体力活动，却需要警醒的管理，使他们也不白吃面包，因为这面包已成为公共财产。也不应看各人在何处隐修院、何处曾将从前所有施予贫困的弟兄。因为所有基督徒组成一个共和国。为此，无论何人在何处向基督徒提供了他们所需之物，那么无论他在何处获得自己所需要的东西，都是从基督的财物中获得的。因为无论他在何处给了这样的人，难道不是基督收下了吗？至于那些在进入这神圣团体之前靠体力劳动谋生的人（进入隐修院的人中大多数属于此类，因为人类中大多数也是如此），他们若不肯劳作，就不该吃饭。因为富人在此基督战斗中谦卑降服，并非为让穷人被抬举至骄傲。确实，在那种生活方式中，议员成为劳动者，而普通劳动者却成为闲散之人，这是极不相称的；并且，当那些曾是房舍和田地的主人们放弃他们的美味而来，普通农夫却变得娇气，这也是极不相称的。\par

34. 然而主说：\Quote{不要为你们的生命忧虑吃什么，也不要为身体忧虑穿什么。}（\BibleRef{玛 6:25}）这很正确，因为祂上面说过：\Quote{你们不能事奉天主而又事奉玛门。}（\BibleRef{玛 6:24}）因为那传福音的人若着眼于可以吃可以穿，他就同时事奉天主（因他传福音）又事奉玛门（因他为了这些需要而传福音）；主说这是不可能的。由此，凡为这些缘故传福音的人，就被证明他不是事奉天主，而是事奉玛门；然而天主可能用他（他自己不知道如何）来促进他人的益处。因为祂在这句话之后紧接着说：\Quote{为此，我告诉你们：不要为你们的生命忧虑吃什么，也不要为你们的身体忧虑穿什么；}（\BibleRef{玛 6:25}）这不是说他们不应该以诚实的方式取得为急需足够之物，而是说他们不应该着眼于此，并为了这些而做他们在传福音时所受命要做的一切。祂称那做事的意向为\Quote{眼目}；上文稍前祂有意要逐步谈到这一点，于是说：\Quote{身体的光就是你的眼睛；如果你的眼睛单纯，你的全身就充满光明；但如果你的眼睛邪恶，你的全身就充满黑暗；}（\BibleRef{玛 6:22-23}）这就是说，你的行为将如何，取决于你做它们的意向如何。因为祂为了谈到这一点，之前已就施舍颁了诫命，说：\Quote{不要在地上为自己积蓄财宝，那里有锈和蛀虫腐蚀，也有贼挖洞偷窃；但要在天上为自己积蓄财宝，那里没有蛀虫和锈腐蚀，也没有贼挖洞偷窃。因为你的财宝在哪里，你的心也必在哪里。}（\BibleRef{玛 6:19-21}）随后祂接着说：\Quote{身体的光就是你的眼睛；}（\BibleRef{玛 6:22}）那就是，那些施舍的人，做施舍时不可怀着取悦于人，或寻求在地上得回报的意向。因此，宗徒在训示弟茂德警告富人时说：\Quote{让他们乐意施舍、分享，为自己积蓄良好的根基，以备将来，好能抓住真正的生命。}（\BibleRef{弟前 6:17-19}）既然主已将那施舍之人的眼目引向来世和天上的赏报，为使他们当眼目单纯时，行为本身充满光明（因为那最后的报应，祂在别处所说的话正是指此：\Quote{谁接待你们，就是接待我；谁接待我，就是接待那派遣我来的。谁因先知的名接待先知，将得到先知的赏报；谁因义人的名接待义人，将得到义人的赏报；谁若只给这些小子中的一个一杯凉水喝，因他是门徒，我实在告诉你们，他决不会失掉他的赏报。}（\BibleRef{玛 10:40-42}）），免得祂在斥责了那些供应必需品给贫穷先知、义人和主门徒之人的眼目之后，那些接受这些事物之人的眼目反而变得败坏，以致为了接受这些事物而愿意像士兵一样事奉基督。祂说：\Quote{没有人能事奉两个主人。}（\BibleRef{玛 6:24}）稍后祂说：\Quote{你们不能事奉天主而又事奉玛门。}（\BibleRef{玛 6:24}）立刻祂加上：\Quote{因此我告诉你们：不要为你们的生命忧虑吃什么，也不要为你们的身体忧虑穿什么。}（\BibleRef{玛 6:25}）\par

35. 至于随后关于天空的飞鸟和田间的百合所说的话，祂说这些话的用意，是叫人不要以为天主不顾祂仆人的需要；因为祂至智的天主眷顾在创造和治理这些事物时，也及于它们。因为不可认为，那亲手劳作的人也不是由祂喂养和穿衣。但为免他们将基督的服役转去谋取这些事物，主以此警告祂的仆人：在这本应归于祂圣事的事奉中，我们不应为这些事物思虑，而应为祂的国和祂的义德思虑；这一切都要加给我们——无论我们是亲手劳作，或因身体软弱不能劳作，或因我们服役本身的这种职务所束缚，以致不能做其他事。因为，主曾说：\BibleQuote{你在患难之日呼求我，我必拯救你，你也要荣耀我}，这并不意味着宗徒不该逃走，不该被人用筐子从城墙上缒下以逃避追捕者的手，而应等待被擒，好让主像从火中拯救三个青年一样拯救他。同样，主也不应说这话：\BibleQuote{若有人在这城迫害你们，就逃往另一城}，因为祂曾说：\BibleQuote{你们若因我的名向父求什么，祂必赐给你们}。因此，如果有人向正在逃避迫害的基督门徒提出这类质问：为何他们不站着，藉呼求天主而透过祂的奇事得到如此拯救，就如同\Person{达尼尔}脱离狮口、\Person{伯多禄}挣脱锁链一样？他们必回答：他们不应试探天主，但若祂愿意，当祂使他们无能为力时，祂便会为他们行同样的事；但是，当祂将逃跑置于他们能力之内时，虽然他们因此得救，但他们仍非由祂拯救不可。同样，对于那些有天主仆人的时间和力量，按照宗徒的榜样和训诫亲手谋生的人，如果有人从\WorkTitle{福音}中就天空的飞鸟（它们不播种、不收割、也不积蓄在仓里）和田间的百合（它们不劳作、也不纺线）提出疑问，他们将轻易回答：若我们也因任何软弱或职务而不能工作，祂会喂养和穿衣我们，如同祂喂养和穿衣那些不做此类工作的飞鸟和百合一样；但是，当我们有能力时，我们不应试探我们的天主；因为这能力本身，我们是从祂的恩赐得来的，而靠此生活，我们是靠祂的慷慨生活，因为祂慷慨地赏赐我们拥有这能力。因此，关于这些必需品，我们并不忧虑；因为当我们能做这些事时，那喂养和穿衣人类的天主喂养和穿衣我们；但是，当我们不能做这些事时，那喂养飞鸟和穿衣百合的天主喂养和穿衣我们，因为我们比它们更贵重。因此，在这我们的服役中，我们也不为明天忧虑；因为我们并非为这些暂时的事物（明天属于它们），而是为那些永恒的事物（在那里永远是今天）而向祂证明自己，使我们不被世俗事务缠身，而能令祂喜悦。\par

36. 既然这些事情如此，圣洁的弟兄，请容许我片刻（因为主藉着你赐给我极大的胆量）来对这些我们的儿子弟兄们说话，我知道你们是如何以同样的爱与他们一同忍受产痛，直到宗徒的训导在他们内成形。天主之仆，基督之兵，你们就是这样掩饰我们最狡猾的仇敌的阴谋吗？他惧怕你们的好名声，即基督的馨香，免得善良的灵魂说：\BibleQuote{我们要追随你香膏的馨香}，从而逃脱他的网罗；他千方百计想用自己的臭气掩盖它，便到处散布了许多穿着隐修士外衣的假善人，在各省游荡，无处被派遣，无处固定，无处站立，无处安坐。有些兜售殉道者的遗骸（如果真是殉道者的话）；另一些夸大他们的繸子和经匣；还有一些编造假故事，说他们听说父母或亲戚在某地还活着，因此他们正前往他们那里。他们都索求，都强要，或是他们有利可图的需要的费用，或是他们虚假圣洁的代价。与此同时，无论他们在何处被发现在作恶，或以任何方式变得臭名昭著，在隐修士这个统称之下，你们的目标就被亵渎——这目标如此良善，如此圣洁，我们以基督之名渴望它，如同在其他地区一样，也在全非洲成长并繁荣。那么，你们岂不被敬虔的热忱点燃？你们的心岂不在里面发热，默想中岂不燃起火焰，以致你们应以善行追击这些人的恶行，你们应为他们断绝卑鄙买卖的借口，因为你们的名誉受损，并在软弱者面前放置绊脚石？因此，请怜悯并同情，并向人类表明，你们不是在安逸中寻求现成的生计，而是通过这目标的窄门窄路，寻求天主的国。你们与宗徒有同样的理由，为那寻找借口的人断绝借口，使那些被他们的臭气窒息的人，因你们的馨香而得以苏醒。\par

37. \BibleQuote{我们不是把重担捆起来搁在你们的肩上，自己却连一个指头也不肯动。}你们要查问并承认我们操劳的工作，以及我们当中一些人的身体软弱，还有我们服务的教会中现已形成的习惯——这习惯让我们自己无暇去做我们劝你们做的那些工作。因为虽然我们可以说：\BibleQuote{有谁当兵自备粮饷呢？有谁栽种葡萄园不吃园里的果子呢？有谁牧养羊群不喝羊奶呢？}然而我指着我们的主耶稣（我因祂的名无畏地说这些话）作证，关系到我的个人方便，我宁愿每天在规定的时辰，如同管理良好的隐修院规则所定，亲手做一些工作，剩余的时间用于读经、祈祷，或从事与圣经有关的工作，也不愿去处理那些关于世俗事务的令人极其烦恼的诉讼——这些诉讼我们必须要么通过裁决了结，要么通过调停截断。这些麻烦正是同一位宗徒加在我们身上的（不是凭他自己的意见，而是凭那藉他说话的那一位），然而我们并未读到他自己曾忍受这些；事实上，他的工作性质不允许这样，因为他奔走于宗徒使命之中。他也没有说：\Quote{如果你们有世俗的诉讼，带到我们这里来}或\Quote{指定我们为你们裁判}，而是说：\BibleQuote{那在教会里被认为是微不足道的人，你们就委派他们。}他又说：\BibleQuote{我说这话是要叫你们羞愧。难道你们中间竟没有一个有智慧的人能判断弟兄之间的事，以致弟兄与弟兄对簿公堂，且是在不信的人面前吗？}因此，他愿意那些住在各地的有智慧的、圣洁的信徒——而非那些为福音事务奔走各处的人——负责审理这些事务。故此，没有任何记载说他曾把时间花在这些事上；而我们却无法为自己开脱，即使我们微不足道，因为他愿意即使是这样的人（在缺乏智慧人的情况下）也被委派，而不愿基督徒的事被带到公众法庭上。然而，我们怀着对永生的盼望，靠着主的安慰承担这些劳动，以期结出忍耐的果实。因为我们是祂教会的仆人，尤其是软弱肢体的仆人——不论我们在同一身体中处于哪个肢体。我略过其他无数的教会挂虑，这些或许无人相信，除非亲身体验过。因此，\BibleQuote{我们不是把重担捆起来搁在你们的肩上，自己却连一个指头也不肯动；}因为实际上，如果这样做不损害我们的职务（鉴察人心的那一位看见了！），我们宁愿去做那些我们劝你们做的事，而不是我们被迫去做的这些事。确实，对你们和我们所有人来说，按照我们的地位和职务劳苦时，\BibleQuote{道路在辛劳和困苦中是狭窄的；}然而，\BibleQuote{当我们怀着希望喜乐时}，\BibleQuote{祂的轭是甘饴的，祂的担子是轻省的}——祂已召唤我们进入安息，祂从涕泣之谷先我们而行，祂自己在那里也并非没有痛苦的压迫。如果你们是我们的弟兄，我们的儿子，我们是你们的同仆，或者在基督里是你们的仆人，那么请听从我们的劝诫，承认我们所吩咐的，接受我们所分施的。但如果我们是法利塞人，把重担捆起来搁在你们的肩上，那么我们所说的你们仍要遵行，即使你们不认同我们所行的。然而，我们被你们论断，或被任何人间的法庭论断，那是极小的事。我们对你们的挂虑，那爱德是多么亲切而深厚，让那赐予者鉴察吧——祂赐予了我们那值得被祂双眼察看的东西。总之，你们随意判断我们吧：宗徒保禄在主内命令并恳求你们，\BibleQuote{要安静——即平和而顺从地——工作，并吃自己的饭。}关于他，我想你们不会相信什么恶事，而那位藉他说话的那一位，你们已相信了祂。\par

38. 这些事，我最亲爱的弟兄\Person{奥勒留}，在基督的怀抱中应受尊敬，主藉着你命令我写关于隐修士工作的事，凡祂赐予我的能力，我没有迟延地写了。我首要的挂虑是：免得那些遵守宗徒训诫的善弟兄，被懒惰和不服从的人甚至称为福音的背信者；这样，那些不工作的人至少会毫无疑问地认为那些工作的人比自己更好。但谁能忍受那些顽梗的人，抗拒宗徒最有益的劝诫，不仅不被当作软弱的弟兄容忍，反而被吹捧为更圣洁的人？以致建立在更健全教义上的隐修院，被这两种诱惑所败坏：即游手好闲的放纵和虚假的圣德之名。那么，让其余的弟兄和儿子们——那些习惯于偏袒这样的人，并因无知而辩护这种僭越行为的人——知道，他们自己最需要被纠正，以便那些人也得以纠正，免得他们\BibleQuote{行善而厌倦}。确实，他们迅速而乐意地供应天主仆人所需要的东西，我们不仅不责备他们，反而衷心拥抱他们；只是不要让他们以错误的仁慈，对他们的未来生命造成更大的伤害，甚于对他们现世生命提供的帮助。\par

39. 因为如果人们不按灵魂的欲望赞美罪人，也不称赞行不义的人，罪就较小。然而，还有什么比希望受属下服从，却拒绝服从上司更不义的呢？我说的是\OriginalTerm{宗徒}{Apostle}，不是我们：以至于他们甚至留\OriginalTerm{长发}{long hair}；对于此事，他根本不愿争议，说：\BibleQuote{若有人爱争辩，我们没有这样的习惯，天主的教会也没有。我现在吩咐这事。}这使我们明白，我们不应寻求巧辩，而应听从发令者的权威。请问，人们如此公然违背宗徒的规诫而留长发，又意味着什么？难道是因为必须放假，连理发匠也不工作？或者，因为他们声称效法福音中的飞鸟，就害怕被拔毛，以免不能飞翔？出于对某些留长发的弟兄的尊敬，我羞于再多说这一过错了；在他们身上，除了这一点，我们找到很多、几乎一切值得尊敬之处。但我们越是在基督内爱他们，就越恳切地劝诫他们。我们也不怕他们的谦逊拒绝我们的劝诫，因为我们也希望在他们跌倒或偏离时，得到像他们这样的人的劝诫。因此，我们劝诫这些圣善的人，不要被虚妄之徒的愚蠢诡辩所动摇，也不要在这悖逆之事上效法那些在其他一切方面远不似他们的人。因为那些贩卖\OriginalTerm{伪善}{hypocrisy}的人，害怕剪短的圣洁比长发更廉价；因为看见他们的人会想起我们所读的那些古人，如\Person{撒慕尔}和其他不剪头发的人。他们却不考虑，那预言的\OriginalTerm{帕子}{veil}与福音中所启示的揭开之间有何区别，宗徒关于后者说：\BibleQuote{你们归向基督时，帕子就除去了。}那藉着置于\Person{梅瑟}面容与以色列百姓之间的帕子所预表的，同样也藉着古时圣人们的长发所预表。因为同一位宗徒说，长发也是代替帕子：这些人在他的权威下受到压制。因为他公开说：\BibleQuote{若男人留长发，是他的\OriginalTerm{羞辱}{disgrace}。}他们说：\Quote{这羞辱正是我们为补赎自己的罪而承担的；}他们摆出模拟谦逊的幌子，以便在它的掩护下进行自高自大的买卖。就好像宗徒在教导骄傲，当他说：\BibleQuote{凡男人祈祷或说预言，若蒙着头，就是羞辱自己的头；}以及\BibleQuote{男人不应蒙头，因为他是天主的肖像和光荣。}因此，那说\BibleQuote{不应}的人，也许不知道如何教导谦卑！然而，如果这羞辱——在福音时期本是一种神圣含义，在预言时期也是如此——竟被这些人当作谦卑之事来追求，那么让他们剃发，并用粗麻布蒙头吧。只是那样就没有他们用来吸引人眼目的把戏了，因为\Person{参孙}不是用粗麻布，而是用他的长发蒙住的。\par

40. 此外，他们为了维护自己的一头长发而发明的那一套（如果言语能表达出来的话），是多么痛苦地可笑啊！他们说：\Quote{男人，宗徒禁止留长发；但那为天国自阉的人已不再是男人了。}啊，无与伦比的昏聩！说这话的人用极端不敬的计谋来武装自己以对抗圣经最显明的宣告，坚持一条弯曲的道路，试图引入一种有害的教义，就是不\BibleQuote{那人不随从恶人的计谋，不站罪人的道路，不坐讥讽者的座位，真是有福。}\BibleRef{咏 1:1}因为如果他日夜默思天主的法律，他会在那里找到宗徒保禄本人，他确实宣认了最高的贞洁，说：\BibleQuote{我愿众人如同我一样。}\BibleRef{格前 7:7}但他在如此生活的同时，也如此说话，显示自己是一个男人。因为他说：\BibleQuote{我作孩子的时候，说话像孩子，理解像孩子，思想像孩子；我成了人，就把孩子的事丢弃了。}\BibleRef{格前 13:11}但我何必提及宗徒呢？因为关于我们的主、救主自己，说这些话的人并不知道他们在想什么。因为除了祂，还有谁是经上所说：\BibleQuote{直到我们众人都达到信仰的一致和对天主子圆满的认识，达到成年人的程度，达到基督圆满年龄的尺度；使我们不再作小孩子，为各种教义之风所飘荡，所卷去，随人欺诈的手法和阴谋，陷于错谬的诡计。}\BibleRef{弗 4:13-14}这些人用这种欺诈欺骗无知的人，用这种阴谋和仇敌的诡计，他们自己转来转去，并在旋转中试图使依附于他们的软弱者的思想也（仿佛）跟着他们旋转，以致他们也不知道自己在哪里。因为他们曾听见或读过所写的：\BibleQuote{你们凡在基督内受洗的，就是穿上了基督；不再有犹太人或希腊人，不再有奴隶或自由人，不再有男人或女人。}\BibleRef{迦 3:27-28}但他们不明白，这是指着肉性的私欲偏情说的，因为在内在的人、我们在心灵的新颖中被更新的地方，并不存在这种性别。那么，他们不要因为自己在男性性别上不工作就否认自己是男人。因为结了婚的基督徒也做这工作，当然并不是因为他们与其余非基督徒甚至牲畜所共有的那个东西而成为基督徒。因为那要么是向软弱让步，要么是向必死的繁殖所付的债务；而另一个则是通过忠信的宣认为把握不朽的永生所表示的。那么，关于男人不可蒙头所吩咐的，在身体上确实预表，但在心灵中——那里有天主的肖像和光荣——才得以实现，这些话本身就表明了：\BibleQuote{男人当然不应当蒙头，因为他是天主的肖像和光荣。}\BibleRef{格前 11:7}因为这幅肖像在哪里，祂自己就宣告了，祂说：\BibleQuote{你们不要彼此说谎；脱去旧人和他的作为，穿上新人，这新人是在认识天主上更新的，按创造他的肖像。}\BibleRef{哥 3:9-10}谁能怀疑这种更新发生在心灵中呢？但如果有人怀疑，让他听一个更清楚的说法。因为在另一处，给出同样的劝告，他这样说：\BibleQuote{正如真理在耶稣内，你们应脱去以前生活的旧人，它因欺骗的情欲而腐败；但应在你们心灵的精神上更新，并穿上新人，它是按照天主的肖像被创造的。}\BibleRef{弗 4:21-24}那么，怎么呢？妇女就没有这种心灵更新，没有天主的肖像吗？谁会这样说呢？但她们身体的性别并不表示这一点；因此她们被吩咐要蒙头。她们作为女人所表示的那个部分，就是可以称为私欲偏情的部分，心灵统治着它，而它自己也服从于它的天主，当生活最正确、最有秩序地进行时。因此，在一个单独的人类个体中，心灵和私欲偏情（前者统治，后者被统治；前者为主，后者为臣），在两个人——男人和女人——身上，就按身体的性别被预表出来。宗徒说到这种神圣的含义，他说男人不应蒙头，而女人应当蒙头。因为心灵越勤于克制来自低处的私欲偏情，就越光荣地向更高处前进，直到整个的人，连同这个现在必死的、脆弱的身体，在最后的复活中穿上不朽和不死的，并且死亡被胜利吞没\BibleRef{格前 15:54}。\par

41. 因此，凡不愿行正义之事者，至少当停止教导错误之事。然而我们在这一番话中所责备的，乃是他人；至于那些因违背宗徒训令而蓄发这一单一过错，以致冒犯并扰乱教会的人——因为有些人由于不愿对他们存任何不善之念，被迫曲解宗徒的明言，另一些人则宁可维护圣经的正确理解，而不愿谄媚任何人，于是在软弱与刚强的弟兄之间便产生了最痛苦且危险的纷争——倘若他们知晓这些事，或许也会毫无迟疑地纠正这一点，而在他们身上，我们赞赏并喜爱其余一切。因此，我们并非责备那些人，而是恳求并郑重地以基督的天主性和人性，以及圣神的爱德，恳请他们不要再把这绊脚石放在基督为之而死的软弱者面前，也不要加重我们心中的忧伤与痛苦——当我们想到，邪恶之人能多么轻易地仿效这种恶行以欺骗人类，当他们看见那些在其它诸多美德上值得我们以基督徒之爱所应尽之礼敬的人身上有此行为时。然而，若在此番劝诫之后，或更确切地说，在此番郑重恳求之后，他们仍决意坚持同一行为，我们将别无他法，唯有忧伤哀恸。愿他们知晓这一点；这已足够。若他们是天主的仆人，他们必怀怜悯。若他们无怜悯之心，我也不愿说更严厉的话。因此，所有这一切——或许我在此所言已超出你我二人事工所愿之冗长——若你赞同，请让我们的弟兄和儿子们知晓，你曾为了他们而俯允将这一重担加于我；但若你觉得有需要删改之处，请通过你的圣座之答复告知于我。\par

\SourceRecord{Translated by H. Browne.；Nicene and Post-Nicene Fathers, First Series；Vol. 3.；Edited by Philip Schaff.；Buffalo, NY: Christian Literature Publishing Co.,；1887.；<http://www.newadvent.org/fathers/1314.htm>.}
